\chapter{Detalles tecnológicos}

\section{Lenguajes y frameworks}
\begin{table}[H]
\centering\small
\caption{Tecnologías por capa (versiones tomadas de los pom.xml y de los archivos Compose).}\label{tab:tec}
\begin{tabularx}{\textwidth}{lYY}
\toprule
\textbf{Capa} & \textbf{Tecnología} & \textbf{Para qué} \\
\midrule
Backend & Java 17 (Temurin), Maven (reactor de 8 módulos) & Detección, seguimiento, alertas, gateway y servicios de apoyo \\
Visión & YOLOv8 exportado a ONNX; ONNX Runtime 1.18.0 & Detectar personas y maletas; estimar la pose del dueño (YOLOv8-pose) \\
Tiempo real & Java-WebSocket 1.5.6; servidor HTTP del JDK con pool de 32 hilos & Empujar alertas; servir panel, fotos y API \\
Datos & H2 2.2.224 (por defecto); PostgreSQL 16.6 (JDBC 42.7.4, HikariCP en el servicio de acceso); MongoDB 7.0.16 (controlador 5.2.1); Redis 7.4.2 (Jedis 5.2.0) & Eventos y alertas; permisos; configuración compartida y su historial; presencia y caché \\
Serialización & Gson 2.11.0 & Mensajes JSON del WebSocket y de la API \\
Panel & HTML, CSS y JavaScript sin \emph{build}; pruebas con Node, jsdom y c8 & Tablero de alertas en vivo, compartir ventana, zonas, reporte \\
Métricas & Micrometer 1.13.6 (registro de Prometheus) & Tráfico, latencia, entrega de alertas y dependencias \\
Pruebas & JUnit 5.10.2; JaCoCo 0.8.14; k6 0.57.0 & Unitarias y de integración; cobertura; carga, estrés y pico \\
\bottomrule
\end{tabularx}
\end{table}

\section{Infraestructura y servicios}
\begin{itemize}
  \item \textbf{Nube:} una instancia EC2 de AWS (AWS Academy Learner Lab); el dominio es \cod{safemode-arsw.duckdns.org} (DuckDNS) y el certificado
  TLS lo obtiene y renueva Caddy con Let's Encrypt.
  \item \textbf{Proxy:} Caddy 2.8 como única entrada pública: \cod{/ws} al WebSocket, \cod{/frames} y \cod{/grafana} detrás de \cod{forward_auth}
  contra \cod{/auth/me}, y el resto al panel.
  \item \textbf{Servicios de terceros:} Google Identity Services (OpenID Connect), la API de bots de Telegram y, de forma opcional, el modelo de
  visión de NVIDIA (el modelo por defecto es un Nemotron sin razonamiento, tras retirarse el anterior).
  \item \textbf{Almacenamiento de evidencia:} fotos PNG en un volumen; la base guarda solo la ruta.
\end{itemize}

\section{DevOps y CI/CD}
El pipeline (\cod{.github/workflows/pipeline.yml}) corre en cada \cod{push} y se compone de \emph{reusable workflows}:

\begin{table}[H]
\centering\small
\caption{Etapas del pipeline.}\label{tab:pipeline}
\begin{tabularx}{\textwidth}{lYY}
\toprule
\textbf{Etapa} & \textbf{Workflow} & \textbf{Qué hace} \\
\midrule
Backend & \cod{reusable-backend.yml} & Java 17, \cod{mvn -B verify}: compila, prueba y mide la cobertura con JaCoCo; sube los informes \\
Frontend & \cod{reusable-frontend.yml} & Node 24, \cod{npm ci}, pruebas del panel con cobertura (c8, \cod{--all}); sube \cod{lcov.info} \\
Quality & \cod{reusable-sonar.yml} & Descarga ambos informes y analiza con SonarCloud; falla si el control de calidad no pasa \\
Deploy & job \cod{deploy} & Despliegue por SSH a la EC2 en \cod{develop} y \cod{main}, con una comprobación posterior de que el panel responde \\
\bottomrule
\end{tabularx}
\end{table}

La etapa de calidad se omite con un aviso (sin romper) mientras el equipo no configure el token y las variables de SonarCloud; con todo
configurado, falla si el control no pasa y, como el despliegue depende de ella, un cambio que no cumple no llega al servidor.

\section{Docker y Docker Compose}
Cada módulo tiene su \cod{Dockerfile}; los archivos Compose se combinan por capas (\cod{-f} varias veces) en lugar de duplicarse:

\begin{table}[H]
\centering\small
\caption{Archivos Compose y su papel.}\label{tab:compose}
\begin{tabularx}{\textwidth}{lY}
\toprule
\textbf{Archivo} & \textbf{Papel} \\
\midrule
\cod{docker-compose.yml} & Base: captura de imagen, visión, base H2 en modo servidor y gateway; proyecto \cod{safemode-arsw} \\
\cod{docker-compose.persistencia.yml} & PostgreSQL, MongoDB, Redis, servicios de acceso y colaboración, Prometheus, Grafana, sonda \cod{/ready} y k6 \\
\cod{docker-compose.tracker.yml} & El tracker en un contenedor (despliegue público) \\
\cod{docker-compose.publico.yml} & Caddy como entrada pública y reinicio automático de los servicios \\
\cod{docker-compose.escalado.yml} & Prototipo de escalado: balanceador Caddy y réplicas del gateway (proyecto \cod{safemode-escalado}, aparte del normal) \\
\cod{docker-compose.sonar.yml} & SonarQube local para ver las calificaciones antes del pipeline \\
\bottomrule
\end{tabularx}
\end{table}

\section{Prometheus, Grafana y k6}
Prometheus (5~s de \emph{scrape} y de evaluación de reglas) descubre cada réplica del gateway por DNS y también recoge los servicios de acceso
y colaboración y la sonda \cod{/ready}; Grafana muestra el tablero del equipo y los de persistencia (13.3.1, porque el formato v2 del tablero no
carga en la 11.5); y k6 corre en un contenedor y envía sus métricas a Prometheus por \emph{remote write}. Los scripts de k6
(\cod{performance/k6/safemode.js}), el ejecutor que levanta el prototipo y deja la base en un estado conocido
(\cod{performance/ejecutar.ps1}), el generador de gráficas (\cod{performance/graficar.py}) y las reglas de alerta quedan versionados y son
reproducibles. Las gráficas calculan la latencia del cliente por ventana a partir de los contadores de k6 (porque k6 exporta sus percentiles
acumulados) y se verificaron contra el resumen de k6: misma media, mismas solicitudes, mismas alertas y mismos errores.
